% Page 2 body. Each NativeSlot uses original portrait coordinates in bp.
\def\OrigTop#1{\vtop{\hrule width0pt height0pt\hbox{#1}}}
\def\OrigPair#1#2{\noindent\hspace*{1.45em}\makebox[.49\linewidth][l]{#1}#2\par}
\def\OrigFour#1#2#3#4{\noindent\hspace*{1.45em}%
  \makebox[.236\linewidth][l]{#1}\makebox[.236\linewidth][l]{#2}%
  \makebox[.236\linewidth][l]{#3}#4\par}
% p002_s01
\NativeSlot{42}{22}{%
  \hspace*{1.45em}后来他又步行\(32\)分钟。\(130+130\div2\times32\)算式求的问题是（\hspace{84bp}）。\par
}
% p002_s02
\NativeSlot{68}{70}{%
  8.\quad 学校手工社团买来\(320\)米长的绳子。做手工时，先剪了\(15\)段，每段长\(6\)米，又剪\par
  \hspace*{1.45em}了\(48\)米备用，还剩下（\hspace{25bp}）米绳子。梳理条件信息，应先求（\hspace{80bp}），\par
  \hspace*{1.45em}再求（\hspace{63bp}），最后求（\hspace{63bp}）。列综合算式为（\hspace{75bp}）。\par
}
\NativeSlot{143}{23}{%
  {\heiti 二、选择（将正确答案的序号填在括号里）。}（每空2分，共10分）\par
}
% p002_s03
\NativeSlot{170}{47}{%
  1.\quad 计算\(4\times25\div25\times4\)的结果，正确的是（\hspace{2em}）。\par
  \OrigFour{A. 1}{B. 16}{C. 0}{D. 都不对}%
}
% p002_s04
\NativeSlot{220}{72}{%
  2.\quad 计算\(21\times[950\div(22+28)-12]\)时，运算顺序是（\hspace{2em}）。\par
  \OrigPair{A. 加\(\to\)减\(\to\)乘\(\to\)除}{B. 除\(\to\)加\(\to\)减\(\to\)乘}%
  \OrigPair{C. 加\(\to\)除\(\to\)减\(\to\)乘}{D. 加\(\to\)除\(\to\)乘\(\to\)减}%
}
% p002_s05
\NativeSlot{295}{99}{%
  3.\quad 如果\(90\times35-25\div5\)先算减法，再算除法，最后算乘法，正确的算式是\par
  \hspace*{1.45em}（\hspace{2em}）。\par
  \OrigPair{A. \(90\times(35-25)\div5\)}{B. \(90\times[(35-25)\div5]\)}%
  \OrigPair{C. \(90\times(35-25\div5)\)}{D. \((90\times35-25)\div5\)}%
}
% p002_s06
\NativeSlot{396}{47}{%
  4.\quad 如图，豆豆不小心将作业本弄脏了，遮住了解决问题的部分内容。按照豆豆\par
  \hspace*{1.45em}所列的算式，本题想要解决的问题是（\hspace{2em}）。\par
}
\NativeSlot{446}{99}{%
  \noindent\hbox to \linewidth{%
    \hskip1.45em\begin{minipage}[t]{260bp}\vspace{0pt}%
      A. 物业费账户里还剩多少元\par
      B. 一共需要缴纳多少元物业费\par
      C. 物业费账户里的钱够扣除几个月\par
      D. 物业费账户里剩余的钱够扣除几个月\par
    \end{minipage}\hfill%
    \OrigTop{\begin{tikzpicture}[x=1bp,y=-1bp,inner sep=0pt,outer sep=0pt]
      \draw[black,line width=.5bp] (0,0) rectangle (187,76);
      \node[anchor=north west,font=\footnotesize,text width=174bp,align=left] at (7,4)
        {爸爸往物业费账户里充值了\(945\)元，第一季度扣除了\(405\)元。};
      \filldraw[fill=black!25,draw=black,line width=.4bp]
        (22,47) .. controls (25,43) and (37,43) .. (44,41)
        .. controls (52,39) and (56,36) .. (64,38)
        .. controls (71,39) and (74,43) .. (80,41)
        .. controls (91,37) and (101,43) .. (114,41)
        .. controls (119,41) and (123,42) .. (124,45)
        .. controls (125,48) and (121,50) .. (115,48)
        .. controls (105,47) and (99,51) .. (92,50)
        .. controls (81,48) and (72,53) .. (62,49)
        .. controls (51,47) and (41,52) .. (27,50)
        .. controls (24,50) and (22,49) .. cycle;
      \node[anchor=west,font=\footnotesize] at (5,64)
        {\((945-405)\div(405\div3)=4\)（个）};
    \end{tikzpicture}}%
  }\par
}
% p002_s07
\NativeSlot{547}{44}{%
  5.\quad 爸爸带豆豆去离家\(2000\)米的图书馆看书。父子俩步行速度是\(70\)米/分，他\par
  \hspace*{1.45em}们从家出发\(25\)分钟后，位置大约在（\hspace{2em}）处。\par
}
\NativeSlot{594}{39}{%
  \noindent\makebox[\linewidth][c]{%
    \begin{tikzpicture}[x=1bp,y=1bp,inner sep=1.2bp,font=\small]
      \draw[line width=.48bp] (0,0)--(169,0);
      \foreach \x/\label in {0/{},84/A,105/B,146/C,169/D}{
        \draw[line width=.48bp] (\x,0)--(\x,2.8);
        \node[above] at (\x,2.3) {\(\label\)};
      }
      \node[below] at (0,-2) {豆豆家};
      \node[below] at (169,-2) {图书馆};
    \end{tikzpicture}%
  }\par
}
\NativeSlot{639}{24}{%
  {\heiti 三、计算。}（共26分）\par
}
% p002_s08
\NativeSlot{665}{71}{%
  1.\quad 直接写得数。（每题1分，共8分）\par
  \OrigFour{\(84\div4=\)}{\(23\times3=\)}{\(360\times4=\)}{\(8+36\div4+8=\)}%
  \OrigFour{\(163-45=\)}{\(98\div7=\)}{\(84\div6=\)}{\((28+32)\div5=\)}%
}
